% command to draw a cube, used as illustration in chapter on functions
\newcommand{\drawCube}[4]{
	% first argument: ll position
	% second argument: x shift
	% 3rd argument: # of arguments
	% 4th argument: size (default = 1)
	% define cube coordinates
	\coordinate (A#1) at (#2,    0,     0);
	\coordinate (B#1) at ({#2+#4},0,     0);
	\coordinate (C#1) at ({#2+#4},#4,0);
	\coordinate (D#1) at (#2,    #4,0);
	\coordinate (E#1) at (#2,     0,#4);
	\coordinate (F#1) at ({#2+#4}, 0,#4);
	\coordinate (G#1) at ({#2+#4},#4,#4);
	\coordinate (H#1) at (#2,    #4,#4);

	% draw cube
	\draw[thick] (A#1) -- (B#1) -- (C#1) -- (D#1) -- cycle;
	\draw[thick] (E#1) -- (F#1) -- (G#1) -- (H#1) -- cycle;
	\draw[thick] (A#1) -- (E#1);
	\draw[thick] (B#1) -- (F#1);
	\draw[thick] (C#1) -- (G#1);
	\draw[thick] (D#1) -- (H#1);

	% define lower left corner of doors
	% door with is defined to always be 0.2*size wide and 0.5*size high. We also define that the space between each door is the remaining width,
	% i.e., if we have three doors, then the remaining space is (size-0.2*size*3)/(3+1)

	% calculate how much empty space to leave between each door
	\pgfmathsetmacro\emptyspace{(#4-0.2*#4*#3)/(#3+1)}

	% define amount of space per door
	\pgfmathsetmacro\doorsspace{0.2*#4}

	\foreach \x in {1,...,#3} {
			\pgfmathsetmacro\llCord{\emptyspace*\x+\doorsspace*(\x-1)}
			% lower door coordinate
			\coordinate (llD\x#1) at ({\llCord+#2},0,#4);

			% draw doors
			\fill[blue!50, rounded corners] (llD\x#1) -- ++(0,{#4/2},0) -- ++({0.2*#4},0,0)node[pos=.5, below](in#1\x){\textcolor{blue}{\texttt{x\x}}} -- ++(0,{-#4/2},0);

			% lower door coordinate for value (for later annotation)
			\coordinate (llD\x#1val) at ({\llCord-.1*#4+#2},{-.2*#4},#4); % the coordinate is simple .1*size to the left of the ll corner and .2*size below.
		}

	% coordinate for result door
	\coordinate (llDOut#1) at ({#2+#4},0,{#4*.2});
	\fill[green!50, rounded corners] (llDOut#1) -- ++(0,{#4/2},0) -- ++(0,0,{#4/2})node[pos=0.8,rotate=45, below, align=center](out#1){\textcolor{black}{\tiny\texttt{out#1}\normalsize}} -- ++(0,{-#4/2},0);

	% draw a 2D rectangle around the cube
	\node[fit=(E#1) (C#1), inner sep=4pt, red, dotted, draw, line width=1pt,rounded corners](dottedRect#1) {};
	% annotate the two-d rect
	\node[above=5pt of dottedRect#1.north west, anchor=west, align=left]{\textcolor{red}{\texttt{def a#1(\textcolor{blue}{\foreach\x in {1,...,#3}{x\x, }...})}\normalsize}};
}
